\documentclass[10pt,english]{article}

\usepackage{amsfonts,url}
\usepackage{amsmath}
\usepackage{amssymb}
\usepackage{hyperref}
% Paper setup
\evensidemargin=0in
\oddsidemargin=0in
\textwidth=6.25in
\topmargin=-0.5in
\headheight=0.0in
\headsep=0.5in
\textheight=9.0in
\footskip=0.5in

\input{institution}

\newcommand{\Interior}[1]{\ensuremath{{#1}^{\circ}}}
\newcommand{\Closure}[1]{\ensuremath{\overline{#1}}}
\newcommand{\Complement}[1]{\ensuremath{{#1}^{c}}}

\newcommand{\Expect}{\ensuremath{\mathrm{E}}}
\newcommand{\vecnot}{\underline}
\newcommand{\RealNumbers}{\ensuremath{\mathbb{R}}}
\newcommand{\RationalNumbers}{\mathbb{Q}}
\newcommand{\ComplexNumbers}{\mathbb{C}}
\newcommand{\Real}{\mathrm{Re}}
\newcommand{\Span}{\mathrm{span}}
\newcommand{\Rank}{\mathrm{rank}}
\newcommand{\Nullity}{\mathrm{nullity}}
\newcommand{\Trace}{\mathrm{tr}}
\newcommand{\Diag}{\mathrm{diag}}
\newcommand{\dd}{\mathrm{d}}
\DeclareMathOperator*{\esssup}{ess\,sup}

% Use < , > inner product
\newcommand{\inner}[2]{{\left\langle #1 \mskip2mu , #2 \right\rangle}}
\newcommand{\tinner}[2]{{\langle #1 \mskip1mu , #2 \rangle}}

% Use < | > inner product
%\newcommand{\inner}[2]{{\left\langle #1 \mskip2mu \middle| \mskip2mu #2 \right\rangle}}
%\newcommand{\tinner}[2]{{\langle #1 \mskip1mu | \mskip1mu  #2 \rangle}}



\begin{document}

\title{ECE 586: Vector Space Methods \\ Lecture 17: Frequently Asked Questions}
\author{Henry D. Pfister \\ \univ}
\date{Revised October 18, 2026}

\maketitle

Here we give a list of questions, and their answers, that were submitted by students after watching the flip video or reviewing the lecture slides.

\section{Best Approximation}
\paragraph{How is best approximation in Banach spaces related with orthogonal projection in Hilbert spaces?}

In a Hilbert space, projection onto a closed subspace gives the unique best approximation in the induced norm (slide 3, ``Orthogonal Projection'').
If the subspace has an orthonormal basis, the coefficient formula is on slide 5, ``Orthogonal Projection onto an Orthonormal Set.''
A general Banach space need not have an inner product, so this orthogonality characterization is not available there.


\paragraph{Is a best approximation always unique, or only in Hilbert spaces?}

For a nonempty closed convex subset of a Hilbert space, a best approximation exists and is unique.
For an arbitrary subset, neither conclusion is automatic, even in a Hilbert space.
In a strictly convex normed space, a best approximation from a convex set is unique whenever it exists.
In finite dimensions, nonempty closed sets have nearest points; strict convexity and convexity of the set then give uniqueness.
A general Banach space need not have either property for all closed convex subsets.


\paragraph{On slide 5, ``Orthogonal Projection onto an Orthonormal Set,'' the orthonormal basis is countable. Shouldn't the superscript of sum be some number $M>0$?}

For a finite orthonormal basis with $M$ vectors, the projection formula is a sum from $1$ to $M$.
For a countably infinite complete orthonormal system, use an infinite sum converging in the Hilbert-space norm.
The latter is not a Hamel expansion: its convergence uses the inner product and completeness.


\paragraph{On slide 3, ``Orthogonal Projection,'' can different $v_1,v_2 \in V$ have the same best approximation $w$?}
Yes, here's a simple 2D example -- think of projecting the vectors (1,1) and (1,3) onto the x-axis. In both cases, you get the vector (1,0).

\paragraph{On slide 6, ``What is a Projection? (1),'' does a projection imply $Tv= v$ and $T^2=T$?}
For a projection, $Tv=v$ holds if $v$ \textit{is already in the range of} $T$ but not in general. Also, $T^2 \triangleq T \circ T = T$ for any linear transform that is a projection (i.e., idempotent).

\paragraph{To check that a matrix P is an orthogonal projection, can we take the range of P and the nullspace of P and verify orthogonality between the sets of basis vectors?}
Yes, but an easier way to do this would be to check that $P$ is both Hermitian and idempotent (see Theorem on slide 8).

\paragraph{I can understand the definition of best approximation, but it's difficult to grasp intuitively. Given a vector, you project it onto an orthogonal set and then add the projections together.  Then, you get the best approximation of this vector. It's difficult to picture this and it seems like magic. }

For an orthogonal basis of nonzero vectors $w_i$, the coefficients are $\inner{x}{w_i}/\|w_i\|^2$; for an orthonormal basis the denominator is one.
If $p=P_W x$, the residual $x-p$ is orthogonal to every vector of $W$.
For any $w\in W$, Pythagoras gives $\|x-w\|^2=\|x-p\|^2+\|p-w\|^2$.
The second term is nonnegative, which explains why $p$ is the closest point.


\paragraph{What are the most important applications of vector approximation?}
Principal Component Analysis (a machine-learning technique often used for dimension reduction), Linear Regression, and function approximation.

\paragraph{Is the definition of a projection in Slides 6 and 7 a more general definition which encompasses the projection we have been working with?}
Yes, the term \emph{projection} can be defined in vector spaces without a norm or inner product.
First, we give this definition (Slides 6-7) and then specialize it to the case of \emph{orthogonal projection} (Slide 8) where an inner product is required for the definition.


\paragraph{In the previous lecture we learned about the second-order bound for functions with a Lipschitz gradient. Is there a first-order bound for it?}
When approximating a function, people often refer to the number of terms in the Taylor series as the order of the approximation.
For example, consider various approximations of $\ln (1+x)$:
\begin{itemize}
\item Zeroth-order approximation: $\ln(1+x) = 0 + O(x)$
\item First-order approximation: $\ln(1+x) = 0 + x + O(x^2)$
\item Second-order approximation: $\ln(1+x) = 0 + x - \frac{1}{2}x^2 + O(x^3)$
\end{itemize}
When the error in this approximation is precisely bounded (rather than left as $O(x^k)$), people sometimes refer to this as a $k$-th order bound).
For $f\colon [0,a] \to \mathbb{R}$ defined by $x \mapsto \ln (1+x)$, we have
\begin{itemize}
\item First-order bound: $|\ln(1+x) - 0| \leq L x $ where $L= \sup_{x\in[0,a]} |f'(x)|$ is the Lipschitz constant of $f$ on $[0,a]$
\item Second-order bound: $|\ln(1+x) - 0 - x| \leq \frac{1}{2} L' x^2$ where $L' = \sup_{x\in[0,a]} |f''(x)|$ is the Lipschitz constant of $f'$ on $[0,a]$
\end{itemize}


\end{document}
